\section{Blender}
Blender es un programa de creación 3D de código abierto y libre con herramientas para el 
pipeline completo, incluyendo modelado, animación, simulación, renderizado y edición de video.
\medskip

Al ser un programa de creación 3D, Blender permite la interacción con los elementos de una escena
a través de la interfaz de usuario, nodos de programación visual y una API para su interacción
mediante scripts. 

\subsection{Arquitectura de Blender}
Blender organiza los datos en escenas, las cuales definen los objetos visibles, cámaras, luces y
opciones de renderizado. Un archivo \textit{.blend}, el formato que Blender usa, puede contener
múltiples escenas, y los datos contenidos en una escena pueden ser vinculados a otra escena, 
permitiendo su reúso y organización \parencite{blendersceneintro}.
\medskip

Los objetos, por su parte, son las entidades contenidas en la escena como mallas 3D, esqueletos,
luces, cámaras y curvas. Cada tipo de objeto está compuesto de dos partes:
\begin{description}
    \item[Objeto] Contiene datos de la transformación.
    \item[Datos de objeto] Contiene datos específicos del tipo de objeto. Por ejemplo, un
    objeto de malla contiene datos de la geometría, materiales y grupos de vértices, mientras
    que una cámara en su lugar almacena configuraciones como la longitud focal o el tamaño del
    sensor.
\end{description}\parencite{blenderobjectintro}
\medskip

Finalmente, las colecciones son agrupaciones de objetos dentro de la escena cuyo objetivo principal
es organizar objetos, además de permitir la manipulación conjunta de algunos parámetros de sus 
contenidos. Las colecciones pueden contener otras colecciones, pero todas están contenidas
dentro de la colección raíz de la escena \parencite{blendercollectionintro}.

\subsection{Editor de nodos}
Blender pone a disposición editores de nodos de geometría y de materiales, los cuales actúan como
una interfaz de programación visual para la creación y manipulación procedural de los datos
de geometría y de materiales, respectivamente.
\medskip

El editor de nodos es utilizado para manipular geometría al agregar un modificador de grupo de nodos
al objeto. Los nodos de geometría permiten instanciar, generar y modificar geometría, así como alterar
sus atributos. Uno de sus usos principales es la generación procedural de geometría, lo que permite 
crear mallas y objetos complejos de forma parametrizada \parencite{blendergeometryintro}.
\medskip

El editor de nodos de materiales (o editor de \textit{shaders}), por otra parte, se utiliza para 
modificar el aspecto de un objeto.
De forma similar a los nodos de geometría, los nodos de material se emplean para la generación
procedural de las texturas utilizadas por el material y su parametrización \parencite{blendershaderintro}.
\medskip

Finalmente, Blender ofrece un editor de nodos de composición para manipular el resultado de una
renderización (como imágenes o videos). Su propósito está enfocado en la edición de posproducción
de contenido, por lo que queda fuera del alcance de este trabajo \parencite{blendercompositeintro}.

\subsection{API de Python}
Blender incluye un intérprete de Python, el cual permanece activo de forma continua, para ejecutar scripts de usuario
y otras herramientas del editor \parencite{blenderapioverview}. Adicionalmente, Blender provee una API estructurada como una biblioteca
de Python.

La API de Blender permite utilizar las mismas herramientas de la interfaz de usuario a través 
de un script, además de facilitar la creación de herramientas nuevas y la extensión del motor gráfico
 \parencite{blenderapiquickstart}.